\scanpage{35}
\mspage{\handwritten{I}}
\margin{\handwritten{$\uparrow$}}
\begin{center}
Cohomologie des catégories de Banach

\rule{8em}{0.4pt}

par Max Karoubi
\end{center}

Soit $\handwritten{\mathscr{C}}$ une catégorie additive \handwritten{\textsuperscript{$(*)$}}. Le \uline{groupe de Grothendieck} $K(\handwritten{\mathscr{C}})$ de la catégorie $\handwritten{\mathscr{C}}$ est le quotient du groupe libre engendré par les classes d'isomorphie $\handwritten{\{}G\handwritten{\}}$ d'objets $G$ de $\handwritten{\mathscr{C}}$ par le sous-groupe engendré par les relations $\handwritten{\{}E+F\handwritten{\}}-\handwritten{\{}E\handwritten{\}}-\handwritten{\{}F\handwritten{\}}=0$. Pour des applications topologiques du groupe de Grothendieck ,voir par exemple \textsuperscript{(1)}. Dans cette note on se propose de définir de même ,pour des catégories additives munies d'une certaine structure ,des groupes $K^n(\handwritten{\mathscr{C}})$ , $n\handwritten{\in}\unsure{Z|\mathbb{Z}}$ ,avec $K^0(\handwritten{\mathscr{C}})\handwritten{\approx}K(\handwritten{\mathscr{C}})$ ,jouissant de certaines propriétés ``cohomologiques'' en relation avec les théorèmes de périodicité de Bott .

I \qquad \uline{Définitions}

1.1 \quad Une \uline{catégorie de Banach} (sous-entendu sur $k$ avec $k=\mathbb{R}$ ou $\mathbb{C}$ ) est la donnée d'une catégorie avec objet nul $\handwritten{\mathscr{C}}$ et , quels que soient les objets $M$ et $N$ de $\handwritten{\mathscr{C}}$ ,d'une structure de $k$-espace de Banach sur l'ensemble $\operatorname{Hom}_{\handwritten{\mathscr{C}}}(M,N)$. Cette donnée est astreinte aux \deleted{\unsure{c|C}}conditions suivantes :

a) Quels que soient les objets $M$ ,$N$ et $P$ de $\handwritten{\mathscr{C}}$ ,l'application
\[
\operatorname{Hom}_{\handwritten{\mathscr{C}}}(M,N)\mathbin{\mathrm{x}}\operatorname{Hom}_{\handwritten{\mathscr{C}}}(N,P)\longrightarrow\operatorname{Hom}_{\handwritten{\mathscr{C}}}(M,P)
\]
définie par la composition des morphismes ,est bilinéaire

\noindent\rule{\linewidth}{0.4pt}

\handwritten{\textsuperscript{$(*)$}} Toutes les catégories considérées ici sont des $\handwritten{\mathscr{U}}$-catégories au sens de Grothendi\handwritten{e}ck \textsuperscript{(3)}. Les objets d'une catégorie forment donc en particulier un ensemble .
% archive sheet number: 34

